\documentclass{my_cv}
\usepackage[skins]{tcolorbox}
\usepackage{academicons}
\usepackage{url}
\usepackage{hyperref}

\hypersetup{
    colorlinks=true,
    linkcolor=blue,
    filecolor=magenta,      
    urlcolor=cyan,
    pdftitle={Overleaf Example},
    pdfpagemode=FullScreen,
}

\begin{document}
\begin{multicols}{2}[
    \titletext{Md}
        {Samiullah}
        {Research Officer, Murdoch Children's Research Institute,\\\hspace{12pt} Royal Children Hospital, Melbourne}
        {PhD in IT, Monash University, Australia}
        {samiullahmdsami@gmail.com}
        {\href{https://github.com/MdSamiullah}{/MdSamiullah}}
        {\href{https://tinyurl.com/u92ce6l}{TINYURL [dot] COM / U92CE6L}}
        {\href{https://linkedin.com/in/md-samiullah}{{/in/md-samiullah}}}
        {\href{https://fb.com/sami.cse}{/sami.cse}}
        {Former Assistant Professor,\\\hspace{12pt} Dept. of Comp. Sci. \& Engi., University of Dhaka}
        {+61 472 796 471}
]
\end{multicols}
% add my learnings: 1. Google cloud Generative AI
% add my learnings: 2. Udemy DJango
% add my learnings: 3. Coursera Statistics course 1
% add my learnings: 4. Coursera Statistics course 2 once done
% add my learnings: 5. Coursera Statistics course 3 once done
% add my learnings: 6. Udemy DL once done
% add my learnings: 7. Udemy NLP once done

% \section{\faFileText}{SUMMARY}
% My strong suits are Programming in Java, Python and C++ in both industry and academia. In my PhD, I have developed an OOBN framework for modelling OOBN applications where OOBN classes have inheritance capability. While implementing its prototype version I used APIs of Hugin, jGraphX and ANTLR. As a proof-of-concept case study, I have re-engineered a real-life large scale BN modelling application that includes discussion with experts and their elicitation. The thesis also includes challenges and guidelines for modelling in iOOBN, modelling BNs from the literature, a new and efficient compilation technique for iOOBN, and an automated hierarchy learning technique.
%%% the following is for statistician officer
% My strong suit is solving problems associated with real-world applications by applying state-of-the-art techniques of Machine Learning, Data Mining, Bayesian Artificial Intelligence and Programming in Java, Python and C++. In my PhD, I have developed an OOBN framework for modelling OOBN applications where OOBN classes have inheritance capability. As a postdoctoral research fellow in MCRI, I am developing an AI model consisting of Bayesian Network (BN), biostatistics, and Machine learning techniques to detect vaccine safety signals as early and as accurately as possible. The model includes BNs, automated reaction mapping and grouping, building models through expert elicitations and automated learning from real-life medical and health data, and statistical measures of continuous and sequential surveillance.

%%% The following is for lecturer
% My strong suit is solving problems associated with real-world applications by applying state-of-the-art techniques of Machine Learning (ML), Data Mining, Bayesian Intelligence, Artificial Intelligence (AI) and programming in Java, Python and C++. My long-term goal is to utilise my 12 years of teaching, research and development experience and the experience of working on real-life complex problems like vaccine safety surveillance as a post-doctoral researcher, in order to building and training a community of AI and ML workers to solve complex real-life and state-of-the-art problems.
% I excel in addressing real-world challenges using cutting-edge techniques in Machine Learning (ML), Data Mining, Bayesian Intelligence, Artificial Intelligence (AI), and programming in languages such as Java, Python, and C++. My ultimate objective is to leverage my 15 years of experience in teaching, research, and development, along with my background in tackling intricate real-world issues like vaccine safety surveillance during my post-doctoral research. I aim to establish and mentor a community of AI and ML professionals dedicated to solving advanced and complex real-life and state-of-the-art problems.



% \begin{multicols}{2}
\section{\faPencil}{WORK EXPERIENCE}
\work{Research Officer / Sep 2022 - Present}%
    {Health Analytics Team}%
    {Murdoch Children's Research Institute}
    {BN modeling, Vaccine Safety Signal Detection, Collaboration, Data Processing, Experimentation, Statistical Analysis, Algorithm Design, Publication}
\vspace{-8pt}    
\work{Assistant Professor / Nov 2020 - Sep 2024}%
    {Department of Computer Science \& Engineering}%
    {University of Dhaka, Dhaka, Bangladesh}
    {Research, Lecturing, Assessment, Curriculum, Syllabus, Fund, Collaboration, Content Preparation}
\vspace{-8pt}    
\work{Lecturer / Feb 2013 - Nov 2020}%
    {Department of Computer Science \& Engineering}%
    {University of Dhaka, Dhaka, Bangladesh}
    {Lecturing, Marking, Assessment, Research}
\vspace{-3pt}    
\work{Research Assistant / Feb 2016 - Dec 2019}%
    {Faculty of Information Technology}
    {Monash University, Australia}%
    {Paper writing, Experimentation, Exam Marking, Lecturing, Assessment, Curriculum}
\vspace{-3pt}    
\work{Lecturer / Jun 2011 - Feb 2013}%
    {Department of Computer Science \& Engineering}%
    {Daffodil International University, Dhaka, Bangladesh}
    {Lecturing, Marking, Assessment, Content Preparation}
\vspace{-8pt}    
\work{Lecturer / Feb 2011 - May 2011}%
    {Department of Computer Science \& Engineering}%
    {City University of Bangladesh, Dhaka, Bangladesh}
    {Lecturing, Marking, Assessment, Content Preparation}
\vspace{-8pt}    
\work{Software Engineer / Dec 2010 - Feb 2011}%
    {Therap Services LLC}%
    {US-based Software Company, Dhaka, Bangladesh}
    {Java, PHP, JSP-Servlets, Oracle, Spring, SQL}

% \columnbreak

\section{\faGraduationCap}{EDUCATION}
    
\school{Information Technology - 2020} %
{\textbf{Doctor of Philosophy}, Monash University, Australia} %
{{\em iOOBN: An Object-Oriented Bayesian Network Modelling Framework with Inheritance.}\\ \href{https://tinyurl.com/snpvgpa}{TINYURL [dot] COM / SNPVGPA}}
    
\school{Computer Scie. \& Engi. - 2012} %
{\textbf{MS}, University of Dhaka, Bangladesh} %
{{\em An efficient approach to mine correlated graphs.} \href{https://tinyurl.com/vrx8khn}{TINYURL [dot] COM / VRX8KHN}}

\school{Computer Scie. \& Engi. - 2010} %
{\textbf{BSC (Honours)}, University of Dhaka, Bangladesh} %
{Queue management based congestion control in wireless body sensor network}


\section{\faBinoculars}{Research Experience}
\begin{itemize}[noitemsep]
    \item \textbf{Publications} (count: 36) \hspace{10pt} \href{https://tinyurl.com/rrs885n}{\aidblp \hspace{5pt} DBLP} \hspace{10pt} \href{https://tinyurl.com/u92ce6l}{\aiGoogleScholarSquare \hspace{5pt} SCHOLAR}
        \begin{itemize}[noitemsep]
            \item \textbf{Journals:} 15
            \item \textbf{Conference:} 17
            \item \textbf{Theses:} 2
            \item \textbf{Workshop paper:} 1
            \item \textbf{Technical Report:} 1
        \end{itemize}
\end{itemize}

\section{\faLineChart}{Taught Courses}
\begin{itemize}[noitemsep]
    \item \textbf{Monash University}
        \begin{itemize}[noitemsep]
            \item FIT 1045: Algorithms and Programming Fundamentals in python
            \item FIT 1008: Introduction to Computer Science
            \item FIT 2004: Algorithms and Data structures
            \item FIT 2100: Operating systems
            \item FIT 3080: Intelligent Systems
            \item FIT 3139: Computational Science
            \item FIT 5047: Intelligent Systems (Postgraduate)
            \item FIT 5142: Advanced Data mining (Postgraduate)
            \item FIT 5211: Algorithms and Data Structures (Postgraduate)
            \item FIT 9133: Programming Foundations in Python (Postgraduate)
        \end{itemize}
    \item \textbf{University of Dhaka, NSU, EWU, AUST, ULab}
        \begin{itemize}[noitemsep]
            \item Introduction to Data Science
            \item Machine learning \& Data mining
            \item Compiler Design 
            \item Advanced database management
            \item Software design pattern
            \item Operating Systems
            \item Theory of Automata/ Computations and Languages
            \item Data Structures
            \item Design and Analysis of Algorithms
            \item Database Design and Analysis
        \end{itemize}
\end{itemize}

\section{\faMoney}{Research Grants}

\begin{itemize}[noitemsep]
    \item BDT 1500000/- (AUD 23978.25, the rate of conversion: 62.56 BDT = 1 AUD) 
ICT Innovation fund from the Department of ICT, Bangladesh Govt. for developing a "DrAi: Artificial Intelligence and Pattern Recognition Driven Assistant for Providing Effective Treatment"
    \item AUD 10,000/- as Regional Collaborations Programme COVID-19 Digital Grants from Australian Academy of Science for “Geo-spatial transfer learning based rumour-spreading trend analysis to detect fake news about COVID-19 and on the effects of its vaccines.”
    \item BDT 385000/- (AUD 6154.42, the rate of conversion: 62.56 BDT = 1 AUD) 
University of Dhaka Centennial Research Grant for developing a "Developing Efficient Technique to Detect False Facts Using Knowledge Graph and Bayesian Network Based Models". 
    \item BDT 200000/- (AUD 3196.93, the rate of conversion: 62.56 BDT = 1 AUD)
Bangladesh Ministry of Science and Technology Special Research Grant for developing a "Data Driven Intelligent Assistant for Medical Practitioners"

\end{itemize}

\section{\faPrint}{Publications}
\subsection{Journals (15)}
\begin{itemize}[noitemsep]
    \item M. T. Alam, C. F. Ahmed, \textbf{M Samiullah}, \& C. K. S. Leung (2023). Discovering Interesting Patterns from Hypergraphs. ACM Transactions on Knowledge Discovery from Data. {\em DOI:} \url{https://doi.org/10.1145/3622940}
    \item MS Islam, PC Kar, \textbf{M Samiullah}, CF Ahmed, CKS Leung, Discovering Probabilistically Weighted Sequential Patterns in Uncertain Databases, Applied Intelligence (Springer),  2023. {\em DOI:} \url{https://doi.org/10.1007/s10489-022-03699-7}
    \item M Momtaz, AA Ferdaus, CF Ahmed, {\bf M Samiullah}, Maximal and closed frequent itemsets mining from uncertain database and data stream, International Journal of Data Science, Vol. 4 (3), (pp. 237-259). {\em DOI:} \url{https://doi.org/10.1504/ijds.2019.102792}
    \item S Akther, MR Karim, \textbf{M Samiullah}, CF Ahmed, Mining non-redundant closed flexible periodic patterns, Engineering Applications of Artificial Intelligence (Elsevier), Vol. 69, pp. 1-23,  2018. {\em DOI:} \url{https://doi.org/10.1016/j.engappai.2017.11.005}
    \item AK Chanda, CF Ahmed, \textbf{M Samiullah}, CKS Leung, A new framework for mining weighted periodic patterns in time series databases, Expert System with Applications (Elsevier), Vol. 79, pp. 207-224, 2017. {\em DOI:} \url{https://doi.org/10.1016/j.eswa.2017.02.028}
    \item S Halder, \textbf{M Samiullah}, Y-K Lee, Supergraph based periodic pattern mining in dynamic social networks, Expert System with Applications (Elsevier), Vol. 72, pp. 430-442,  2017. {\em DOI:} \url{https://doi.org/10.1016/j.eswa.2016.10.033}
    \item T Hashem, MR Karim, {\bf M Samiullah}, CF Ahmed, An efficient dynamic superset bit-vector approach for mining frequent closed itemsets and their lattice structure, Expert Systems with Applications, Vol. 67, (pp. 252-271). {\em DOI:} \url{https://doi.org/10.1016/j.eswa.2016.09.023}.
    \item AU Ahmed, CF Ahmed, \textbf{M Samiullah}, N Adnan, CKS Leung, Mining interesting patterns from uncertain databases, Information Sciences (Elsevier), Vol. 354, pp. 60–85,  2016. {\em DOI:} \url{https://doi.org/10.1016/j.ins.2016.03.007}.
    \item \textbf{M Samiullah}, CF Ahmed, A Fariha, MR Islam, N Lachiche, Mining frequent correlated graphs with a new measure, Expert Systems with Applications (Elsevier), Vol. 41, Issue 4, pp. 1847–1863,  2014. {\em DOI:} \url{https://doi.org/10.1016/j.eswa.2013.08.082}.
    \item A Fariha, CF Ahmed, CKS Leung, \textbf{M Samiullah}, S Pervin and L Cao, A New Framework for Mining Frequent Interaction Patterns from Meeting Databases, Engineering Applications of Artificial Intelligence (Elsevier), Vol. 45, pp. 103–118,  2015. {\em DOI:} \url{https://doi.org/10.1016/j.engappai.2015.06.019}.
    \item A K Chanda, S Saha, MA Nishi, \textbf{M Samiullah}, CF Ahmed, An efficient approach to mine flexible periodic patterns in time series databases, Engineering Applications of Artificial Intelligence (Elsevier), Vol. 44, pp. 44–63, 2015. {\em DOI:} \url{https://doi.org/10.1016/j.engappai.2015.04.014}.
    \item T Hashem, CF Ahmed, {\bf M Samiullah}, S Akther, BS Jeong, S Jeon, An efficient approach for mining cross-level closed itemsets and minimal association rules using closed itemset lattices, Expert systems with applications, Vol. 41 (6), (pp. 2914-2938). {\em DOI:} \url{https://doi.org/10.1016/j.eswa.2013.09.052}.
    \item MA Nishi, CF Ahmed, \textbf{M Samiullah}, and BS Jeong, Effective periodic pattern mining in time series databases, Expert Systems with Applications (Elsevier), Vol. 40, Issue 8, pp. 3015–3027,  2013. {\em DOI:} \url{https://doi.org/10.1016/j.eswa.2012.12.017}.
    \item \textbf{M Samiullah}, D. Albrecht, A. E. Nicholson, \& C. F. Ahmed (2021). A Review on Probabilistic Graphical Models and Tools. Dhaka University Journal of Applied Science and Engineering, 6(2), 82-93. {\em DOI:} \url{https://doi.org/10.3329/dujase.v6i2.59223}
    \item M. T. Alam, C. F. Ahmed, \& \textbf{Md Samiullah} (2022). A Vertex-extension based Algorithm for Frequent Pattern Mining from Graph Databases. Dhaka University Journal of Applied Science and Engineering, 7(1), 58-65. {\em DOI:} \url{https://doi.org/10.3329/dujase.v7i1.62887}
\end{itemize}

\subsection{Conferences (17)}
\begin{itemize}[noitemsep]
    \item \textbf{M Samiullah}, Ann Nicholson, David Albrecht (2023). Shareable and Inheritable Incremental Compilation in iOOBN. Pacific Rim International Conference on Artificial Intelligence, (Accepted for publication).
    \item A.M. Aahad, Khondker Bin Yamin, \textbf{M Samiullah}, C. F. Ahmed, C. K. S. Leung, E. W.R. Madill, A. G.M. Pazdor (2023). Connectable and Independent Junction Tree-Based Compilation Technique of Object-Oriented Bayesian Networks. International Conference on Tools with Artificial Intelligence, (Accepted for publication).
    \item W. Islam, R Karim, \textbf{M Samiullah}, C. F. Ahmed, C. K. S. Leung (2023). An Efficient Approach to Learn an Effective Hierarchy of a Set of OOBN Classes. International Conference on Ubiquitous \textbf{I}nformation \textbf{M}anagement and \textbf{Com}munication (IMCOM), (Accepted for publication).
    \item I. Alam, \textbf{M Samiullah}, U. Kabir, S. Woo, C. K. S. Leung (2023). SREMIC: Spatial Relation Extraction-based Malware Image Classification. International Conference on Ubiquitous \textbf{I}nformation \textbf{M}anagement and \textbf{Com}munication (IMCOM), (Accepted for publication).
    \item \textbf{M Samiullah}, Ann Nicholson, David Albrecht (2022). Automated construction of an Object-Oriented Bayesian Network (OOBN) Class Hierarchy. International Conference on Tools with Artificial Intelligence, (pp. ). {\em DOI:} \url{https://doi.org/}.
    \item RR Haque, CF Ahmed, \textbf{M Samiullah}, CK Leung (2021). UFreS: A New Technique for Discovering Frequent Subgraph Patterns in Uncertain Graph Databases. IEEE International Conference on Big Knowledge (ICBK) (pp. 253-260). {\em DOI:} \url{https://doi.org/10.1109/ickg52313.2021.00042}.
    \item M.T. Alam, C.F. Ahmed, \textbf{M. Samiullah}, C.K. Leung (2021). Mining Frequent Patterns from Hypergraph Databases. Pacific-Asia Conference on Knowledge Discovery and Data Mining (pp. 3-15). {\em DOI:} \url{https://doi.org/10.1007/978-3-030-75765-6_1}.
    \item M.T. Alam, C.F. Ahmed, \textbf{M. Samiullah}, C.K. Leung (2021, May). Discriminating frequent pattern based supervised graph embedding for classification. Pacific-Asia Conference on Knowledge Discovery and Data Mining (pp. 16-28). {\em DOI:} \url{https://doi.org/10.1007/978-3-030-75765-6_2}.
    \item \textbf{M Samiullah}, TX Hoang, D Albrecht, A Nicholson, K Korb (2017), iOOBN: A Bayesian Network Modelling Tool Using Object Oriented Bayesian Networks with Inheritance, International Conference on Tools with Artificial Intelligence, (pp. 1218-1225). {\em DOI:} \url{https://doi.org/10.1109/ictai.2017.00185}.
    \item T Alam, SA Zahin, {\bf M Samiullah}, CF Ahmed (2017), An efficient approach for mining frequent subgraphs, International Conference on Pattern Recognition and Machine Intelligence. (pp. 486-492). {\em DOI:} \url{https://doi.org/10.1007/978-3-319-69900-4_62}.
    \item MBUZ Shajib, {\bf M Samiullah}, CF Ahmed, CK Leung, AGM Pazdor (2016), An efficient approach for mining frequent patterns over uncertain data streams, International Conference on Tools with Artificial Intelligence, (pp. 980-984). {\em DOI:} \url{https://doi.org/10.1109/ictai.2016.0151}.
    \item CF Ahmed, {\bf M Samiullah}, N Lachiche, M Kull, P Flach (2015), Reframing in frequent pattern mining, International Conference on Tools with Artificial Intelligence, (pp. 799-806). {\em DOI:} \url{https://doi.org/10.1109/ictai.2015.118}.
    \item R Sultana, D Showkat, \textbf{M Samiullah}, and AR Chowdhury, Reconstructing Gene Regulatory Network with Enhanced Particle Swarm Optimization, ICONIP 2014, Part II, LNCS 8835, pp. 229–236, 2014. {\em DOI:} \url{https://doi.org/10.1007/978-3-319-12640-1_28}.
    \item {\bf M Samiullah}, CF Ahmed, MA Nishi, A Fariha, SM Abdullah, M Islam (2013), Correlation mining in graph databases with a new measure, Asia-Pacific Web Conference, (pp. 88-95). {\em DOI:} \url{https://doi.org/10.1007/978-3-642-37401-2_11}.
    \item S Halder, {\bf M Samiullah}, AMJ Sarkar, YK Lee (2012), Movie swarm: Information mining technique for movie recommendation system, International Conference on Electrical and Computer Engineering, (pp. 462-465). {\em DOI:} \url{https://doi.org/10.1109/icece.2012.6471587}.
    \item {\bf M Samiullah}, SM Abdullah, AFMIH Bappi, S Anwar (2012), Queue management based congestion control in wireless body sensor network, International Conference on Informatics, Electronics \& Vision, (pp. 493-496). {\em DOI:} \url{https://doi.org/10.1109/iciev.2012.6317349}.
    \item {\bf Md Samiullah}, PC Kar, MS Islam, MT Alam \& CF Ahmed (2023), Dr. AI: A Heterogeneous Clinical Decision Support System for Personalized Health Care, International Conference on Information and Communication Technology. {\em DOI:} \url{https://doi.org/10.1007/978-981-19-2394-4_29}
\end{itemize}

\section{\faBook}{Thesis}
\begin{itemize}[noitemsep]
    \item \textbf{PhD Thesis:} {\em iOOBN: An Object-Oriented Bayesian Network Modelling Framework with Inheritance.}\\ \href{https://tinyurl.com/snpvgpa}{TINYURL DOT COM / snpvgpa}
    \item \textbf{MS Thesis:} {\em An efficient approach to mine correlated graphs.} \href{https://tinyurl.com/vrx8khn}{TINYURL DOT COM / vrx8khn}
\end{itemize}

\section{\faGroup}{ROLES}
\begin{itemize}[noitemsep]
    \item \textbf{Conference Reviewer} 
        \begin{itemize}[noitemsep]
            \item Pacific Asia Knowledge Discovery in Databases
            \item International Conference on Machine Learning
            \item AAAI
        \end{itemize}
    \item \textbf{Journal Reviewer} 
    \begin{itemize}[noitemsep]
            \item IEEE Transactions on Systems, Man and Cybernetics: Systems
            \item Applied Intelligence
            \item International Journal of Data Science and Analytics
            \item Dhaka University Journal of Applied Science and Engineering
            \item IEEE Access
        \end{itemize}
    \item \textbf{Senior Academic Staff} 2018-2019, For units FIT1045, FIT9133, Faculty of IT, Monash University
    \item \textbf{Syllabus and Curriculum} development committee member: FIT1045, FIT9133, Faculty of IT, Monash University
    \item \textbf{Academic Integrity Committee} member: FIT1045, FIT9133, Faculty of IT, Monash University
    \item \textbf{Syllabus Design} committee member: Machine Learning Track, Department of computer Science and Engineering, University of Dhaka
    \item \textbf{Adjunct faculty}: North South University, Independent University of Bangladesh, East West University, Ahsanullah University of Science \& Technology, Green Univ. of Bangladesh
\end{itemize}

\section{\faBuilding}{Other Experiences}
\begin{itemize}[noitemsep]
    \item \textbf{Project Coordination Team member:} University of Dhaka AIS (Accounting Information Systems)
    \item \textbf{Project Administration Team member:} National University of Bangladesh Online Admission Management System
    \item \textbf{Consultant and expert team member:} Titas Gas (National Gas Transmission and Distribution Company of Bangladesh) Hardware and Software Testing
    \item \textbf{Project Coordination, development and maintenance Team member:} Central Admission Management System of the University of Dhaka
\end{itemize}

\section{\faSoccerBallO}{Professional Training}
{\bf Murdoch Children's Research Institute (MCRI):} 

MCRI provided the following professional development training as a Research Officer:

\begin{itemize}[noitemsep]
    \item COVID - 19 Safe Training
    \item RCH Emergency Response curricula
    \item MCRI and VCGS Core Compliance Training
    \item Research Data Management
    \item Performance and Development - Team Member
\end{itemize}

\noindent\\
{\bf Monash University:} 

Monash University provided the following professional development training as a PhD student and teaching associates:

\begin{itemize}[noitemsep]
    \item Research Strategies and Skills
    \item Advanced Research Methods
    \item Research Integrity
    % \item Research Integrity: Social and Behavioural Sciences
    % \item Research Integrity: Engineering and Technology
    \item Ethics and Professional Conduct
    \item Understanding Mental Health
    \item Data Protection and Privacy
    \item Introduction to Cyber Security
    \item Equal Opportunity Online Training
\end{itemize}

\noindent\\
{\bf University of Dhaka:} 

University of Dhaka provided the opportunity to following professional development training as a faculty member:

\begin{itemize}[noitemsep]
    \item BDITEC (Bangladesh IT-engineers Examination Center) Management
    \item ITEE (Information Technology Engineers Examination) Training of Trainers
\end{itemize}

\noindent\\
{\bf Online Learning Platforms:} 

In Udemy and Coursera, I have completed the following courses as a part of my learning, self-development, and standing out of the crowd journey:

\begin{itemize}[noitemsep]
    \item Udemy: Complete Data Science Bootcamp
    \item Coursera: Understanding and Visualizing Data with Python
    \item Udemy: Python and Django Full Stack Web Developer Bootcamp
    \item Udemy: React - The Complete Guide
    \item Udemy: Deep Learning A-Z
    \item Udemy: NLP - Natural Language Processing with Python
\end{itemize}

\section{\faBuilding}{BN Modelling Experience}

\begin{itemize}[noitemsep]
    \item Re-engineered Western Grassland Reserve project (Melbourne) as a case-study
    \item Developed various BNs from the literature for experimentation
    \item Developing a BN model for Vaccine Safety Signal detection
    \item BN construction through expert elicitation for vaccine safety signal detection
\end{itemize}

\section{\faList}{SKILLS}

\textbf{BN Software APIs:} Hugin (Java, Python), BN Learn (python), Netica (Java, Python)

\noindent\textbf{Languages:} Python, Java, PHP, C++, C, JavaScript, HTML, CSS, Bash, SQL, R

\noindent\textbf{Tools:} PyTorch, NumPy, Pandas, dJango, React, Power-BI, Spring-Hibernate, jGraphX, PEP, ANTLR, Latex, LEX-YACC, Weka, Hugin-Expert, Netica, Android Studio, Mendeley, Endnote

\noindent\textbf{Other:} Version Control (Git and Git-Hub), Web, Data science and ML, Data Mining.

\section{\faBuilding}{Development Experience}
\begin{itemize}[noitemsep]
    \item Developing the Trustworthy-AI tool for modelling Bayesian network in Medical domain using React framework
    \item Developing iOOBN framework for OO-Bayesian Network Modelling (Java) using Hugin (API), jGraphX (API) and ANTLR (API)
    \item Full-stack web developer (PHP, JS, JQuery, HTML, CSS, Bootstrap)
Online Admission and Fee Payment System, Central Admission Office, University of Dhaka
    \item Chief Investigator and Project Manager, Smart prescription generator, (Flask:Python) \href{http://draibd.com}{DRAIBD [dot] COM}
    \item Project Coordinator, Freelancer-Client network (Node js and React js), \href{https://somadha.com}{SOMADHA [dot] COM}
    \item Project Coordinator, Tele-medicine App, (Node js and React js), \href{https://sobardakter.com}{SOBARDAKTER [dot] COM}
    \item Implementing MaxSPRT (A statistical meaure for sequential surveillance) for vaccine safety signal detection under SAEFVIC data and settings.
    \item Automated Mapping of Reactions to MEDRA-Code.
    \item Developing an automated system for Unit Marking Break-down Summary (UMBS) in Monash University whil working as a TA during my PhD.
\end{itemize}

\section{\faTrophy}{AWARDS}
\begin{itemize}[noitemsep]
    \item Postgraduate Publication Award 2019, FIT, Monash University
    \item Travel grant from ABNMS to attend ABNMS 2018 Conference in Adelaide, Australia
    \item Travel grant to attend ICTAI 2017 in Boston, USA
    \item Australian Postgraduate Award: Monash University
    \item International Postgraduate Research Scholarship (IPRS): Monash University 
    \item Merit Scholarship (in B.Sc.) in the University of Dhaka (Held in 2009) 
\end{itemize}

\section{\faSoccerBallO}{INTERESTS}

\begin{itemize}[noitemsep]
    \item \textbf{Research:} Machine Learning, Data Mining, Bayesian Artificial Intelligence, Bayesian Statistics, Probabilistic Analysis 
    \item Software development, UI/UX Design
\end{itemize}

% \end{multicols}

\noindent\rule{18.75cm}{0.6pt}
\vspace{20pt}

\begin{multicols}{3}

\section{\faUniversity}{Reference 1}
    
\work{Professor Ann E Nicholson}%
    {Dean}%
    {Faculty of Information Technology, Monash University, Australia\\ 
    \faPhone \hspace{5pt} +61 4 480 19439 \hspace{20pt}\\
    \faEnvelope \hspace{5pt} ann.nicholson@monash.edu}
    {PhD Supervisor, Postdoc Supervisor}



\columnbreak

\section{\faInstitution}{Reference 2}

\work{Ingrid Zukerman}%
    {Professor}%
    {Faculty of Information Technology, Monash University, Australia\\ 
    \faPhone \hspace{5pt} +61 3 9905 5202 \hspace{20pt}\\
    \faEnvelope \hspace{5pt} ingrid.zukerman@monash.edu}
    {Supervisor of Units tutored}


\columnbreak

\section{\faInstitution}{Reference 3}

\work{Carson Kai-Sang Leung}%
    {Professor}%
    {Department of Computer Science, Univ. of Manitoba, Canada\\
    \faPhone \hspace{5pt}  +1 (204) 474-8677 \hspace{20pt}\\
    \faEnvelope \hspace{5pt} Carson.Leung@umanitoba.ca}
    {Research Collaborator, Mentor}

\end{multicols}

\end{document}
